\chapter{Soon, Speedily, Or In Speed?}

\section{The Options For \textit{En Tachei}}

The opening does not only say that God shows his servants certain events. It says that God shows them \textgreek{ἃ δεῖ γενέσθαι ἐν τάχει} (\textit{ha dei genesthai en tachei}, \textit{the things that must happen in speed}). The question is what kind of claim the final phrase makes.

Many English translations render \textit{en tachei} as ``soon'' or ``shortly.'' That reading is understandable. If something must happen \textit{en tachei}, it may be natural to hear temporal nearness: the events are about to occur, or at least are near from the perspective of the speaker and first hearers. This is the reading that has often shaped the larger interpretation of the book of Revelation. Once the opening says ``soon,'' the reader must ask whether the events were fulfilled near John's time, whether divine time differs from human time, whether imminence can stretch across centuries, or whether the phrase creates an unresolved tension.

A second option is ``quickly'' or ``speedily.'' This reading does not focus first on how soon the events begin, but on the manner or speed with which they happen. On this view, the phrase may suggest rapid occurrence when the appointed action takes place. The events are not necessarily dated by the opening; they are characterized by speed.

A third option is the deliberately wooden rendering used earlier in this book: ``in speed.'' This is not polished English. It is a control gloss. Its purpose is to prevent the interpreter from deciding too quickly between temporal nearness and modal speed. It keeps the phrase open long enough for the structure to test it.

Each option assumes something about the events. ``Soon'' assumes that the phrase primarily locates the events in relation to the hearers' time. ``Speedily'' assumes that the phrase primarily describes how the events occur. ``In speed'' assumes less at the start. It says only that speed belongs to the things that must happen, while leaving the precise force to be decided by the passage's own system.

The book's method requires the third approach as the starting point. That does not mean ``soon'' is impossible. Greek grammar permits a temporal reading, and the wider book contains nearness language that cannot be ignored. But possibility is not governance. The question is not whether \textit{en tachei} can be rendered ``soon.'' The question is whether the opening structure requires ``soon'' as the controlling meaning.

\section{Timing And Manner}

The difference between timing and manner may appear small, but it changes the interpretive frame.

If \textit{en tachei} means ``soon,'' the opening makes a direct claim about temporal nearness. The things shown to the servants must happen near the time of the disclosure. That reading pushes the interpreter toward questions of fulfillment. Which events happened soon? How soon is soon? Does the phrase refer to the beginning of the events, the whole complex of events, or the decisive end? If the events did not all happen near John's time, how should the statement be defended?

Those questions are legitimate if the phrase is temporal. The problem is that they can take control of the book before the opening structure has finished its work. The reader may begin with chronology and then force the transmission sequence to serve that chronological question. The servants become an audience waiting for near events. The disclosure becomes a timetable. The book becomes a problem of delay.

If \textit{en tachei} means ``speedily,'' the focus shifts. The things that must happen are characterized by speed, not necessarily by immediate onset. The phrase then speaks less like a date marker and more like a description of event-movement. When the necessary things unfold, they unfold in speed. This reading fits the force of the noun behind the phrase better than a purely calendrical paraphrase does, but it can also be overstated. If ``speedily'' is used to remove all temporal pressure, it becomes another shortcut.

The opening does not let the interpreter choose timing against manner as though only one can matter. Speed language can carry temporal force without functioning as a simple date. A fast event is temporal; it happens in time. But it is not the same thing as a near event. The issue is whether \textit{en tachei} points mainly to proximity before the events begin or to the speed/urgency that belongs to their happening.

The verb phrase also matters. The words are not merely ``things in speed.'' They are \textgreek{δεῖ γενέσθαι} (\textit{dei genesthai}, \textit{must happen}). The phrase combines necessity and occurrence. These are not optional possibilities. They are things that must come to pass. \textit{En tachei} modifies that field of necessary happening. It is attached to the events as events, not to the act of writing the book, not to John's testimony as testimony, and not directly to the reader's act of hearing.

That observation narrows the question. The opening is not first saying, ``This book must be read soon.'' It is not first saying, ``John must testify quickly.'' It is not even first saying, ``The disclosure is sent quickly.'' It says that the things shown and signified to the servants must happen \textit{en tachei}. The timing or speed belongs to the content of the showing and signifying.

This is why ``soon'' is too narrow as the book's working default. It turns a phrase about necessary happening into a single question about chronological nearness. The phrase may include nearness, but the structure asks for a broader category: speed as a feature of divinely necessary events disclosed to servants.

\section{Testing The Phrase In The Structure}

The expanded form gives the test. The phrase occurs in the C line and is supplied to the C' line:

\begin{center}
\begin{longtable}{@{}p{0.1\textwidth}p{0.78\textwidth}@{}}
\toprule
Pair & Expanded line \\
\midrule
C & \textit{deixai tois doulois autou ha dei genesthai en tachei} \\
C' & \textit{kai [tois doulois autou ha dei genesthai en tachei] esemanen} \\
\bottomrule
\end{longtable}
\end{center}

The things that must happen in speed are both shown and signified. That placement is decisive for method. \textit{En tachei} is not an isolated timing slogan at the front of the book. It is part of the content transferred through paired acts of divine communication.

This explains why the servant chapter had to come first. The recipients are not detached spectators. They are God's servants. The events are disclosed to those who stand under command, loyalty, and obedience. If the things must happen in speed, the speed is not given to satisfy curiosity about a schedule. It belongs to the way servants are prepared to receive, recognize, and obey within the disclosed order of events.

The repetition near the end of the book confirms the importance of this phrase without removing the need for local analysis. Revelation 22:6 again speaks of what must happen \textit{en tachei}. That repetition frames the book, but it does not by itself decide whether the phrase means ``soon'' in a simple chronological sense. It shows that the opening language remains programmatic. The same necessary happenings disclosed at the beginning are still in view at the end.

On this analysis, the best working rendering is still ``in speed,'' with ``speedily'' as the smoother English option when a normal sentence is needed. ``Soon'' should not control the interpretation unless the local argument requires temporal nearness. The opening itself does not require that reduction. It requires that the things shown and signified to God's servants be understood as divinely necessary happenings marked by speed.

This conclusion does not solve every timing question in the book of Revelation. It does not decide how every later vision relates to John's historical moment, the churches of Asia, later history, or final consummation. It does something more limited and more important for this study: it prevents the opening from being reduced to an imminence claim before its communication system has been understood.

The result also sets up the next question. If the things that must happen in speed are both shown and signified, then the mode of communication matters. The issue is not only when the events happen or how rapidly they unfold. The issue is how God communicates such events to servants. That requires returning to the paired verbs themselves: showing and signifying.
